% Exercitium_Vol1.tex
% Exercitium (curso 2013-14)
% José A. Alonso Jiménez <https://jaalonso.github.io>
% Sevilla, 14 de Octubre de 2025
% ======================================================================

\documentclass[a4paper,12pt,twoside]{book}

%%%%%%%%%%%%%%%%%%%%%%%%%%%%%%%%%%%%%%%%%%%%%%%%%%%%%%%%%%%%%%%%%%%%%%%%%%%%%%
%% § Paquetes adicionales                                                   %%
%%%%%%%%%%%%%%%%%%%%%%%%%%%%%%%%%%%%%%%%%%%%%%%%%%%%%%%%%%%%%%%%%%%%%%%%%%%%%%

% Configuración para XeLaTeX
\usepackage{tocloft}
\setlength{\cftbeforechapskip}{2ex}
\setlength{\cftbeforesecskip}{0.5ex}
\setlength{\cftchapnumwidth}{2em}
\setlength{\cftsecnumwidth}{14mm}
\cftsetindents{chapter}{0em}{\cftchapnumwidth}

\usepackage{fontspec}
\usepackage{xunicode}
\usepackage{xltxtra}
\defaultfontfeatures{Ligatures=TeX,Numbers=OldStyle} % inicial
\setromanfont{DejaVu Sans} % inicial
\setmonofont{DejaVu Sans Mono}[Scale={0.90}] % inicial

% \setsansfont{Arial}

% \setmainfont{DejaVu Math TeX Gyre}
% \setmonofont{DejaVu Math TeX Gyre}[Scale={1.1}]

% Notas: La lista de fuentes disponibles se obtiene con fc-list

\usepackage[spanish]{babel}          % Castellanización.
\usepackage{a4wide}
\usepackage{amssymb, amsmath, mathtools}

\usepackage{newunicodechar}
\newunicodechar{⟹}{\ensuremath{\mathnormal\Longrightarrow}}
\usepackage[dvipsnames]{xcolor}

\usepackage{minted}
\usemintedstyle{emacs}
\definecolor{bg}{rgb}{0.95,0.95,0.95}
\newminted{lean}{bgcolor=bg, autogobble=true}

\setminted{fontsize=\small,baselinestretch=1}

\newcommand{\ejercicio}[1]{\inputminted[firstline=7]{haskell}{../src/#1.hs}}

\linespread{1.05}                  % Distancia entre líneas
\setlength{\parindent}{2em}        % Indentación de comienzo de párrafo
% \deactivatetilden                  % Elima uso de ~ para la eñe
\raggedbottom                      % No ajusta los espacios verticales.

\usepackage[%
  colorlinks=true,
  urlcolor=blue,
  % pdftex,
  pdfauthor={José A. Alonso <jalonso@us.es>},%
  pdftitle={Exercitium Vol. 1},%
  pdfstartview=FitH,%
  bookmarks=false]{hyperref}

\setcounter{tocdepth}{1}
\setcounter{secnumdepth}{4}

% No indenta los item
\usepackage{enumitem}
\setlist[itemize]{leftmargin=*}

%%%%%%%%%%%%%%%%%%%%%%%%%%%%%%%%%%%%%%%%%%%%%%%%%%%%%%%%%%%%%%%%%%%%%%%%%%%%%%
%% § Cabeceras                                                              %%
%%%%%%%%%%%%%%%%%%%%%%%%%%%%%%%%%%%%%%%%%%%%%%%%%%%%%%%%%%%%%%%%%%%%%%%%%%%%%%

\usepackage{fancyhdr}

\addtolength{\headheight}{\baselineskip}

\pagestyle{fancy}

\cfoot{}

%\fancyhead{}
%\fancyhead[RE]{\mdseries\sffamily \nouppercase{\leftmark}}
%\fancyhead[LO]{\mdseries\sffamily \nouppercase{\rightmark}}
%\fancyhead[LE,RO]{\mdseries\sffamily \thepage}

% \fancyhead{}
% \fancyhead[RE]{\mdseries \nouppercase{\leftmark}}
% \fancyhead[LO]{\mdseries \nouppercase{\rightmark}}
% \fancyhead[LE,RO]{\mdseries \thepage}

\fancyhead{}
\fancyhead[RE,LO]{\mdseries \nouppercase{\leftmark}}  % Capítulo en ambos lados
\fancyhead[LE,RO]{\mdseries \thepage}                 % Página en esquinas

%%%%%%%%%%%%%%%%%%%%%%%%%%%%%%%%%%%%%%%%%%%%%%%%%%%%%%%%%%%%%%%%%%%%%%%%%%%%%%
%% § Título                                                                 %%
%%%%%%%%%%%%%%%%%%%%%%%%%%%%%%%%%%%%%%%%%%%%%%%%%%%%%%%%%%%%%%%%%%%%%%%%%%%%%%

\title{
  {\LARGE Exercitium \\
    {\Large Ejercicios de programación funcional con Haskell \\
      {\normalsize (Volumen 1: curso 2013--14)
    }}}}
\author{
  \href{https://jaalonso.github.io/}{José A. Alonso Jiménez}}
\date{\vfill \hrule \vspace*{2mm}
  \begin{tabular}{l}
      Grupo de Lógica Computacional \\
      Dpto. de Ciencias de la Computación e Inteligencia Artificial \\
      Universidad de Sevilla  \\
    Sevilla, 14 de octubre de 2025
  \end{tabular}\hfill\mbox{}}

%%%%%%%%%%%%%%%%%%%%%%%%%%%%%%%%%%%%%%%%%%%%%%%%%%%%%%%%%%%%%%%%%%%%%%%%%%%%%%
%% § Documento                                                              %%
%%%%%%%%%%%%%%%%%%%%%%%%%%%%%%%%%%%%%%%%%%%%%%%%%%%%%%%%%%%%%%%%%%%%%%%%%%%%%%

% \includeonly{}

% \includexmp{licencia}

\begin{document}

\maketitle
\newpage

\noindent
Esta obra está bajo una licencia Reconocimiento--NoComercial--CompartirIgual
2.5 Spain de Creative Commons. 

\vspace*{4ex}

\begin{center}
\fbox{
\begin{minipage}{\textwidth}
\vspace*{1ex}
\noindent                     
{\bf Se permite:}
\begin{itemize}                         
\item copiar, distribuir y comunicar públicamente la obra
\item hacer obras derivadas
\end{itemize}
         
\noindent       
{\bf Bajo las condiciones siguientes:}

\begin{minipage}[t]{20mm}
  \vspace*{0mm}\hspace*{1em}
  \includegraphics[scale=0.35]{deed.png}
\end{minipage}%
\begin{minipage}[t]{14cm}
  \vspace*{0mm}
  {\bf Reconocimiento}. Debe reconocer los créditos de la  obra de la manera
  especificada por el autor.  
\end{minipage}
\\[1ex]

\begin{minipage}[t]{20mm}
  \vspace*{0mm}\hspace*{1em}
  \includegraphics[scale=0.35]{deed-eu.png}
\end{minipage}%
\begin{minipage}[t]{14cm}
  \vspace*{2mm}
  {\bf No comercial}. No puede utilizar esta obra para fines comerciales.
\end{minipage}
\\[1ex]

\begin{minipage}[t]{20mm}
  \vspace*{0mm}\hspace*{1em}
  \includegraphics[scale=0.35]{deed_002.png}
\end{minipage}%
\begin{minipage}[t]{14cm}
  \vspace*{-2mm}
  {\bf Compartir bajo la misma licencia}. Si altera o transforma esta obra, o
  genera una obra derivada, sólo puede distribuir la obra generada bajo una
  licencia idéntica a ésta. 
\end{minipage}

\begin{itemize}    
\item Al reutilizar o distribuir la obra, tiene que dejar bien claro los
  términos de la licencia de esta obra. 
\item Alguna de estas condiciones puede no aplicarse si se obtiene el permiso
 del titular de los derechos de autor.
\end{itemize}
\vspace*{1ex}
\end{minipage}}
\end{center}

\vspace*{4ex}

Esto es un resumen del texto legal (la licencia completa). Para ver una copia
de esta licencia, visite 
\href{http://creativecommons.org/licenses/by-nc-sa/2.5/es/}
     {\url{http://creativecommons.org/licenses/by-nc-sa/2.5/es/}}
o envie una carta a Creative Commons, 559 Nathan Abbott Way, Stanford,
California 94305, USA.

%%% Local Variables: 
%%% mode: latex
%%% TeX-master: t
%%% End: 

\newpage

\tableofcontents
\clearpage

\renewcommand{\chaptername}{Ejercicio}

\chapter*{Introducción}

% \mbox{} \hspace*{1cm}

\begin{quote}
  ``\textit{The chief goal of my work as an educator and author is to
  help people learn to write beautiful programs.}''

  (Donald Knuth en
  \href{http://www.paulgraham.com/knuth.html}{Computer programming as an art})
\end{quote}

\vspace* {1cm}

Este libro es una recopilación de las soluciones de los ejercicios
propuestos en el blog
\href{https://jaalonso.github.io/exercitium}
     {Exercitium}\
     \footnote{\url{https://jaalonso.github.io/exercitium}}
durante el curso 2013--14.

El principal objetivo de Exercitium es servir de complemento a la
asignatura de
\href{https://jaalonso.github.io/cursos/i1m-13}
     {Informática}\
     \footnote{\url{https://jaalonso.github.io/cursos/i1m-13}}
de 1º del Grado en Matemáticas de la Universidad de Sevilla.

Con los problemas de Exercitium, a diferencias de los de las
\href{https://web.archive.org/web/20250614171615/https://www.cs.us.es/~jalonso/cursos/i1m-13/ejercicios/ejercicios-I1M-2013.pdf}
     {relaciones}\
     \footnote{\url{https://web.archive.org/web/20250614171615/https://www.cs.us.es/~jalonso/cursos/i1m-13/ejercicios/ejercicios-I1M-2013.pdf}},
se pretende practicar con los conocimientos adquiridos durante todo el
curso, mientras que con las relaciones están orientadas a los nuevos
conocimientos.

En cada ejercicio se muestra distintas soluciones, se comprueba con
QuickCheck su equivalencia y se compara su eficiencia.

El código de los ejercicios del libro se encuentra en
\href{https://github.com/jaalonso/Exercitium}
     {GitHub}\
     \footnote{\url{https://github.com/jaalonso/Exercitium}}

% \chapter{Ejercicios de Abril de 2014}

\chapter{Iguales al siguiente}
\ejercicio{Iguales_al_siguiente}

\chapter{Ordenación por el máximo}
\ejercicio{Ordenados_por_maximo}

\chapter{La bandera tricolor}
\ejercicio{Bandera_tricolor}

\chapter{Determinación de los elementos minimales}
\ejercicio{ElementosMinimales}

\chapter{Mastermind}
\ejercicio{Mastermind}

\chapter{Primos consecutivos con media capicúa}
\ejercicio{Primos_consecutivos_con_media_capicua}

\chapter{Anagramas}
\ejercicio{Anagramas}

\chapter{Primos equidistantes}
\ejercicio{Primos_equidistantes}

% \chapter{Ejercicios de Mayo de 2014}

\chapter{Suma si todos los valores son justos}
\ejercicio{Suma_si_todos_justos}

\chapter{Matriz de Toeplitz}
\ejercicio{Matriz_Toeplitz}

\chapter{Máximos locales}
\ejercicio{Maximos_locales}

\chapter{Lista cuadrada}
\ejercicio{Lista_cuadrada}

\chapter{Segmentos maximales con elementos consecutivos}
\ejercicio{Segmentos_consecutivos}

\chapter{Valores de polinomios representados con vectores}
\ejercicio{Valor_de_un_polinomio}

\chapter{Ramas de un árbol}
\ejercicio{Ramas_de_un_arbol}

\chapter{Alfabeto comenzando en un carácter}
\ejercicio{Alfabeto_desde}

\chapter{Numeración de ternas}
\ejercicio{Numeracion_de_ternas}

\chapter{Ordenación de estructuras}
\ejercicio{Ordenacion_de_estructuras}

\chapter{Emparejamiento binario}
\ejercicio{Emparejamiento_binario}

\chapter{Amplia columnas}
\ejercicio{Amplia_columnas}

\chapter{Regiones determinadas por n rectas del plano}
\ejercicio{Regiones}

\chapter{Elemento más repetido de manera consecutiva}
\ejercicio{Mas_repetido}

\chapter{Número de pares de elementos adyacentes iguales en una matriz}
\ejercicio{Pares_adyacentes_iguales}

\chapter{Mayor producto de las ramas de un árbol}
\ejercicio{Mayor_producto_de_las_ramas_de_un_arbol}

\chapter{Biparticiones de una lista}
\ejercicio{Biparticiones_de_una_lista}

\chapter{Trenzado de listas}
\ejercicio{Trenzado_de_listas}

\chapter{Números triangulares con n cifras distintas}
\ejercicio{Triangulares_con_cifras}

\chapter{Enumeración de árboles binarios}
\ejercicio{Enumera_arbol}

\chapter{Algún vecino menor}
\ejercicio{Algun_vecino_menor}

\chapter{Reiteración de una función}
\ejercicio{Reiteracion_de_funciones}

% \chapter{Ejercicios de Junio de 2014}

\chapter{Pim, Pam, Pum y divisibilidad}
\ejercicio{PimPamPum}

\chapter{Código de las alergias}
\ejercicio{Alergias}

\chapter{Índices de valores verdaderos}
\ejercicio{Indices_verdaderos}

\chapter{Descomposiciones triangulares}
\ejercicio{Descomposiciones_triangulares}

\chapter{Número de inversiones}
\ejercicio{Numero_de_inversiones}

\chapter{Separación por posición}
\ejercicio{Separacion_por_posicion}

\chapter{Emparejamiento de árboles}
\ejercicio{Emparejamiento_de_arboles}

\chapter{Eliminación de las ocurrencias aisladas}
\ejercicio{Elimina_aisladas}

\chapter{Ordenada cíclicamente}
\ejercicio{Ordenada_ciclicamente}

\chapter{Órbita prima}
\ejercicio{Orbita_prima}

\chapter{Divisores de un número con final dado}
\ejercicio{Divisores_con_final}

\chapter{Descomposiciones de x como sumas de n sumandos de una lista}
\ejercicio{Descomposiciones_con_n_sumandos}

\chapter{Selección hasta el primero que falla inclusive}
\ejercicio{Seleccion_con_fallo}

\chapter{Buscaminas}
\ejercicio{Buscaminas}

\chapter{Mayor sucesión del problema 3n+1}
\ejercicio{Mayor_sucesion_3n_mas_1}

\chapter{Filtro booleano}
\ejercicio{Filtro_booleano}

\chapter{Entero positivo de la cadena}
\ejercicio{Entero_positivo_de_la_cadena}

\chapter{N gramas}
\ejercicio{N_gramas}

\chapter{Sopa de letras}
\ejercicio{Sopa_de_letras}

\chapter{Intercalación de n copias}
\ejercicio{Intercala_n_copias}

\chapter{Eliminación de n elementos}
\ejercicio{Elimina_n_elementos}

% \chapter{Ejercicios de Julio de 2014}

\chapter{Límite de sucesiones}
\ejercicio{Limites_de_sucesiones}

\chapter{Empiezan con mayúscula}
\ejercicio{Empiezan_con_mayuscula}

\chapter{Renombramiento de un árbol}
\ejercicio{Renombra_arbol}

\chapter{Divide si todos son múltiplos}
\ejercicio{Divide_si_todos_multiplos}

\chapter{Ventana deslizante}
\ejercicio{Ventana_deslizante}

\chapter{Representación de Zeckendorf}
\ejercicio{Representacion_de_Zeckendorf}

\chapter{Elemento más cercano que cumple una propiedad}
\ejercicio{Codigo_Morse}

\chapter{Producto cartesiano de una familia de conjuntos}
\ejercicio{Producto_cartesiano}

\chapter{Todas tienen par}
\ejercicio{Todas_tienen_par}

\chapter{Sucesiones pucelanas}
\ejercicio{Sucesiones_pucelanas}

\chapter{Producto de matrices como listas de listas}
\ejercicio{Producto_de_matrices_como_listas_de_listas}

\chapter{Inserción en árboles binarios de búsqueda}
\ejercicio{Insercion_en_arboles_binarios_de_busqueda}

\chapter{Matriz permutación}
\ejercicio{Matriz_permutacion}

\chapter{Números con todos sus dígitos primos}
\ejercicio{Numeros_con_digitos_primos}

\chapter{Cadenas de ceros y unos}
\ejercicio{Cadenas0y1}

\chapter{Clausura de un conjunto respecto de una función}
\ejercicio{Clausura}

\chapter{Sustitución en una expresión aritmética}
\ejercicio{Sustitucion_en_una_expresion_aritmetica}

\chapter{Laberinto numérico}
\ejercicio{Laberinto_numerico}

\end{document}

%%% Local Variables:
%%% mode: latex
%%% TeX-master: t
%%% End:
